\section{Markov Chains, Hidden Markov Models, and the Markov Blanket}
\label{sec:markov-chain-blanket}

\subsection{Markov Chains}

A Markov Chain is a graphical abstraction of a stochastic process that
transitions between a finite number of states, used to model systems whose
future evolution depends only on their present state and not on the full
history that led to it. This memorylessness property is formalised as the
\textbf{Markov property}: for a discrete-time process $\{X_t, t \ge 0\}$
taking values in a countable state space,
\[
  \Pr[X_{t+1} = j \mid X_t = i, X_{t-1} = k, \ldots, X_0 = m]
  = \Pr[X_{t+1} = j \mid X_t = i] = p_{ij}, \qquad \forall t \ge 0.
\]
A finite-state Markov chain with $N$ states can be represented either as a
graph (a node per state, a directed weighted arc for every non-zero
transition probability $p_{ij}$) or, equivalently, as an $N \times N$
row-stochastic transition matrix $[M] = \{p_{ij}\}$ with $\sum_j p_{ij} = 1$
for every $i$. This duality is useful in practice: the graphical view makes
the qualitative structure of the process easy to inspect, while the matrix
view makes it straightforward to apply standard linear-algebra machinery
(e.g.\ computing the stationary distribution as the matrix's leading
eigenvector).

\subsection{Hidden Markov Models}

A \textbf{Hidden Markov Model (HMM)} extends a Markov chain to the case where
the states themselves are not directly observed. Formally, an HMM is a
5-tuple $(Q, \Sigma, \Pi, A, B)$, where $Q = \{q_1, \ldots, q_N\}$ is a finite
set of hidden states, $\Sigma = \{s_1, \ldots, s_M\}$ is a finite set of
observable symbols, $\Pi = \{\pi_i\}$ is the initial-state probability
vector, $A = \{a_{ij}\}$ is the hidden-state transition matrix, and
$B = \{b_i(v_k)\}$ is the emission-probability matrix giving the probability
of observing symbol $v_k$ while in hidden state $i$. The transition between
hidden states is itself assumed to follow a first-order Markov chain as
defined above; only the emissions generated by each hidden state are
observed directly.

Three canonical problems are associated with an HMM $\lambda = (\Pi, A, B)$,
each with a classical dynamic-programming solution:
\begin{description}
  \item[Evaluation] Compute the probability $p(O \mid \lambda)$ of an
    observed emission sequence $O = o_1 o_2 \cdots o_T$ under the model —
    solved by the \textbf{Forward} and \textbf{Backward} algorithms.
  \item[Decoding] Find the most likely hidden-state path
    $q^* = \arg\max_{q} p(q, O \mid \lambda)$ that produced an observed
    sequence — solved by the \textbf{Viterbi} algorithm.
  \item[Training] Find the parameters $\lambda^* = \arg\max_{\lambda}
    p(O \mid \lambda)$ that best explain an observed sequence — solved by
    the \textbf{Baum-Welch} algorithm (an instance of Expectation-Maximisation
    that repeatedly calls the Forward and Backward algorithms to re-estimate
    $\Pi$, $A$, and $B$) \autocite{rabiner1986hmm}.
\end{description}
In this project, the HMM's hidden states correspond to latent market regimes (e.g.\ bull/bear/stagnant), while the emissions are the observed values (or discretised bins) of each financial/macro/geopolitical
series.

\subsection{The Markov Blanket}

The \textbf{Markov blanket} of a node $v$ in a DAG is defined as the union of
$v$'s parents, $v$'s children, and the other parents of $v$'s children (its
children's co-parents) \autocite{pearl1988probabilistic}. Its defining
property is that it is a minimal sufficient summary of everything else in the network relevant to $v$: every other node in the graph is conditionally independent of $v$ once the Markov blanket of $v$ is known,
i.e.\
\[
  X_v \perp\!\!\!\perp X_u \mid B(X_v) \qquad \text{for every } u \notin B(X_v) \cup \{v\}.
\]
This property has two direct, practical consequences for this project.
First, it gives risk management a precise, minimal answer to "what do I need
to monitor to explain the oil-price node?": once the price node's Markov
blanket is known, every other variable in the 24-series panel becomes
redundant for predicting it. Second, and more immediately, it is the
computational engine behind the constraint-based Inferred Causality
structure-learning algorithm developed next (Section
\ref{sec:inferred-causality}): learning each node's Markov blanket first lets
the algorithm discard the overwhelming majority of candidate conditioning
sets before testing for undirected arcs, which is what makes constraint-based
learning tractable at all on a panel of this size
\autocite{pellet2008markov}.
